\subsection{Part�cula Livre (Solu��o Mais Simples)}

A solu��o mais simples para a part�cula livre\index{part�cula livre} de Dirac, � obtida
quando o vetor momento, $\versor{p}$, apontar na dire��o positiva do eixo $z$, nesse caso,
a equa��o de Dirac, j� em sua forma matricial, toma a seguinte apar�ncia:
%
%
\[
\begin{pmatrix}
mc^2 & 0      & pc  & 0 \\
0    & mc^2   & 0     & - pc \\
pc & 0      & mc^2  & 0 \\
0    & - pc & 0     & mc^2
\end{pmatrix}
\begin{pmatrix}\psi_1 \\ \psi_2 \\ \psi_3 \\ \psi_4 \end{pmatrix} =
\epsilon\begin{pmatrix}\psi_1 \\ \psi_2 \\ \psi_3 \\ \psi_4 \end{pmatrix}
\]

Repare que por simplicidade de nota��o, foi definido que, $p_3 = p$ .De forma que se tenha
o seguinte conjunto de equa��es homog�neas nas vari�veis, $\psi_i$, sendo, $i = 1,2,3,4$:
%
%
\[\begin{array}{ccccccccc}
(\epsilon - mc^2)\psi_1 &+& 0\psi_2 &-& pc\psi_3 &-& 0\psi_4    &=& 0\\
0\psi_1 &+& (\epsilon - mc^2)\psi_2 &-& 0\psi_3 &+& pc\psi_4  &=& 0\\
- pc\psi_1 &-& \psi_2 &+& (\epsilon - mc^2)\psi3 &+& 0\psi_4 &=& 0\\
0\psi_1 &+& pc\psi_2 &+& 0\psi_3 &+& (\epsilon - mc^2)\psi_4 &=& 0
\end{array}
\]

O determinante\index{determinante} (polin�mio caracter�stico) possui como resultado
(ra�zes), os seguintes
valores para o auto valor, $\epsilon$:
%
%
\[\epsilon_+ = \sqrt{mc^2 + (pc)^2} \quad , \quad \epsilon_- = - \sqrt{mc^2 + (pc)^2}\]

Para o auto valor\index{auto valor}, $\epsilon_+$, existem duas solu��es linearmente
independentes.
%
%
\[\psi_1^{(1)} = 1
\ \ , \ \
\psi_2^{(1)} = 0
\ \ , \ \
\psi_3^{(1)} = \frac{pc}{\epsilon_+ + mc^2}
\ \ , \ \
\psi_4^{(1)} = 0
\]
%
%
\[
\psi_1^{(2)} = 0
\ \ , \ \
\psi_2^{(2)} = 1
\ \ , \ \
\psi_3^{(2)} = 0
\ \ , \ \
\psi_4^{(2)} = - \frac{pc}{\epsilon_+ + mc^2}
\]

Da mesma forma, temos que, para o auto valor, $\epsilon_-$, existem duas solu��es
linearmente independentes.
%
%
\[\psi_1^{(1)} = \frac{pc}{\epsilon_- - mc^2}
\ \ , \ \
\psi_2^{(1)} = 0
\ \ , \ \
\psi_3^{(1)} = 1
\ \ , \ \
\psi_4^{(1)} = 0
\]
%
%
\[\psi_1^{(2)} = 0
\ \ , \ \
\psi_2^{(2)} = - \frac{pc}{\epsilon_- - mc^2}
\ \ , \ \
\psi_3^{(2)} = 0
\ \ , \ \
\psi_4^{(2)} = 1
\]
